\section{Materials and methods}
\noindent\textit{Draft status, 2 November 2026. This section distinguishes documented procedures and counts from the protocol for remaining work. Corpus-wide full-text eligibility and a final included-study count have not been established.}

\subsection{Review design and analytical scope}
The review combines systematic retrieval with a comparative, evidence-linked synthesis of three relationships: building material assessments and uncertainty; domain-specific digital twins; and digital twins and uncertainty mechanisms. The analytical scope extends from observation and inference of building attributes through material inventories, temporal flows and recovery pathways to environmental assessment and decision support. S1--S5 were used as organising concepts in the searches. Their correspondence with the structures of integrated assessment frameworks is subject to explicit comparison; they are not claimed to have emerged independently of the search design or to represent universal software layers.

Publications constitute the unit of retrieval and screening. The unit of synthesis is a documented transformation or interface within a publication. One publication may contribute to several stages or evidence streams, but is counted once in a DOI-deduplicated publication registry. Separate publications are not assumed to represent independent studies until their underlying study families have been assessed.

\subsection{Search families and information sources}
Archived searches cover Scopus, Web of Science Core Collection and IEEE Xplore. The stage-specific DT searches were executed on 1 November 2026 and the stage-specific DT-plus-uncertainty searches on 2 November 2026. Earlier S1--S3 tables require reconciliation with their original search logs before an overall search interval can be reported. Executable strings, query identifiers, export counts and native fields are retained in the source directories.

Four stage-oriented retrieval families are distinguished: Base, Base plus uncertainty (U), Base plus DT, and Base plus DT plus U. General DT and DT-plus-U searches provide a separate interdisciplinary resource. These labels describe retrieval families, not mutually exclusive eligibility classes. The historical DT-stage queries broadened some building contexts and observation terms; consequently they are not literal nested subsets of the original Base searches. They are reported separately and must not be used to infer prevalence ratios without query harmonisation. The expanded search for integrated frameworks has no imposed publication-year or publisher restriction; its identification and screening log remains to be completed and is not incorporated into the archived database counts.

Scopus used title, abstract and keyword fields; Web of Science used Topic searches, including Keywords Plus; IEEE used the recorded search scope and native export format. These field scopes are not identical. The DT-stage protocol records no intentional year, language, document-type or open-access filter. Restrictions in earlier search versions must be taken from their own logs. A displayed result total is distinguished from a successfully exported record set. In particular, general DT annual facets do not constitute a complete record-level export.

\subsection{Record management, screening and retrieval}
Native exports were preserved and bibliographic fields harmonised while retaining source and query membership. Exact normalised DOIs were used for conservative merging; native database identifiers preserved identity where DOIs were absent. Similar titles and conflicting DOIs were flagged for adjudication rather than automatically equated. Duplicate decisions and bibliographic corrections were retained in audit files. Deduplication status differs across families and is reported explicitly.

Title and abstract screening distinguished direct material-chain relevance, complementary methods, secondary evidence and unrelated records. Operational-only energy studies, generic project management and isolated material tests without the required assessment connection were distinguished from core material-flow evidence. An uncertainty or DT keyword alone did not establish eligibility. Incomplete information was retained as unresolved where a full text was necessary. Background and transferable-method records were retained separately rather than represented as core inclusions.

The S4--S5 recheck documents AI-assisted reading and recorded rationales for all 9,587 canonical records, including a second review of 946 previously unclear cases. Multiple assistant readers and consistency checks were used; this is not independent human dual screening. Automated proposal stages and later semantic reassessment are distinct in the audit trail. A complete full-text eligibility assessment of the resulting candidate set has not been performed. Elicit was explored but is not credited with screening the full corpus. Tool identities, versions and any unavailable execution details require a consolidated reporting log before submission.

PRISMA 2020 terminology distinguishes records, reports sought, reports obtained, reports assessed and included studies (Page et al., 2021, DOI: 10.1136/bmj.n71). The accompanying flow snapshots report only reconstructable transitions. Downloading a PDF is not counted as reading it, and an archived appraisal is not treated as final inclusion. Failure to retrieve a report is distinct from exclusion after assessment.

\subsection{Digital twin definitions and data extraction}
Definitions are treated as evidence to be compared rather than resolved by the presence of the term digital twin. The extraction protocol separately records the authors' definition and cited definition source, the physical referent, digital representation, data connection, update mechanism and demonstrated analytical or decision function. This operational description draws on the characterisation and definition literature already cited in the introduction; it is not asserted to be a universally accepted threshold.

Manual versus automated exchange, one-way versus two-way exchange, periodic versus event-driven updates, and prediction versus feedback control are recorded separately. A BIM, GIS or three-dimensional model is not automatically coded as an operational twin. Conversely, absence of a reported feature is coded as not reported, not as proof that the feature is absent. Conceptual architectures, prototypes and operational deployments are distinguished. Domain studies that provide useful information mechanisms without satisfying an operational DT description remain identifiable as complementary evidence.

For uncertainty, the protocol records the affected quantity, source, representation, quantification method, dependence assumptions, transformation and downstream consequence. Spatial variability, epistemic uncertainty and alternative scenarios are not collapsed into one probabilistic category. Interface coding distinguishes documented transfer, documented representation change or loss, and insufficient reporting. Any claim of reduced uncertainty requires a stated comparison basis. Spatial extent and grain, temporal validity and horizon, and stakeholder information requirements are recorded as contextual conditions. Each substantive mechanism claim should be linked to a page, section, equation or figure; missing locators remain explicit. These fields extend existing appraisals and are not yet complete across all retrieved records.

\subsection{Comparative analysis and evidence synthesis}
Existing relational analyses aggregate coded working appraisals into stage--uncertainty, stage--DT-function and DT--uncertainty-mechanism relationships. The denominator is stated for each analysis. Multiple memberships are preserved, and counts across stages or streams are not summed as independent publications. Frequencies indicate representation in the assessed set, not effectiveness or absence from the wider field. Retrieval-only records are excluded from full-text mechanism claims.

The intended final synthesis links an observed interface finding to an information requirement, a reported DT capability and its conditions of transfer. Directly demonstrated relationships are distinguished from cross-study interpretations and research proposals. A conceptual architecture is an organising result of the review, not an implemented or validated system. Literature silence is separated from a documented limitation. No meta-analysis or pooled effectiveness estimate is planned because the current evidence combines heterogeneous methods, contexts and study types.

\subsection{Current evidence status and completion requirements}
The S4--S5 combined Base/U audit reconciles 19,963 query occurrences to 14,869 database-distinct records and 9,587 adjudicated canonical records. Screening retained 2,932 clear full-text candidates, 327 unresolved candidates and 351 secondary-evidence candidates; 5,977 records were outside that candidate set. Of the clear candidates, 392 belonged to the uncertainty-query subset. These are not final included studies.

The broad Base-plus-DT retrieval contains 4,124 occurrences and 2,447 conservatively unique records, with 42 possible title-duplicate groups unresolved. Base-plus-DT-plus-U contains 903 occurrences and 534 conservatively unique records, with seven title groups unresolved. Their corpus-wide screening and full-text dispositions cannot yet be reported. The focused DT bridge search is a separate retrieval: 83 occurrences yielded 43 records, of which 14 were selected for full-text work and 12 have archived full-text assessments. It is not represented as a screened subset of the broad DT retrieval without a record-level provenance join.

The working synthesis registry contains 52 distinct publications across 53 stream memberships. Forty-four have archived full-text-based appraisals: 20 in S--U, 12 in S--DT and 13 in DT--U, with one shared assessed publication. Eight registry records lack an archived full-text appraisal. These figures are readiness counts, not a complete PRISMA inclusion endpoint. Final reporting requires reconciliation of the early Base records, completion of the broad DT-family screening, full-text eligibility decisions, the supplemental search log, study-family checks and documented human verification. No prospective registration or independent dual-human review is claimed.
